\chapter{Limitaciones}
\label{sec:paper-limites}

Un reporte interno responsable declara lo que \emph{no} funciona
con tanta claridad como lo que funciona. Este capítulo lista las
limitaciones conocidas, agrupadas por ámbito.

\section{Arquitectura}

\subsection*{Dos aristas residuales core$\to$platform}

El modelo shearing-layers exige que \code{core} no dependa de
\code{platform}. Quedan dos violaciones:

\begin{itemize}
    \item \code{core/graph-engine} depende de
    \code{platform/sources} para el adapter que hidrata el grafo
    desde una fuente canonical. Lo correcto sería invertir: un
    trait \code{CanonicalSource} en \code{core} y la
    implementación en \code{platform}.
    \item \code{core/graph-engine} depende de \code{platform/dsl}
    para la representación estructural de una hipótesis. Idem:
    un AST mínimo en \code{core} y el parser en \code{platform}.
\end{itemize}

Ambas fueron conservadas durante el rediseño por costo de cambio
relativo bajo y alcance creciente del plan. Están listadas en
\filepath{docs/_archive/phase2/FOLLOWUPS.md} como follow-ups
F2.11 y F2.12.

\section{Modelo ML}

\subsection*{GNN no entrenado para FortiAnalyzer}

El clasificador v2 fue entrenado sobre Sentinel AAD. Datasets de
firewall (FortiAnalyzer), IAM cloud (Cognito) y SaaS (Zoho)
carecen de modelo propio. El comportamiento es correcto (retorna
0 con peso 0, no degrada), pero significa que GraphHunter hoy no
aporta el quinto componente de scoring en la mayoría de
deploys. Training pipeline reproducible queda como trabajo
futuro.

\subsection*{Calibración de pesos}

Los pesos por defecto
$(0.25, 0.25, 0.30, 0.20, 0.00)$ son una elección informada por
inspección manual sobre datasets sintéticos. No hay evaluación
cuantitativa contra un set etiquetado grande. Un analista con
contexto organizacional puede ajustar, pero no hay método
automático de búsqueda de pesos óptimos.

\section{Tests}

\subsection*{Flakes conocidos}

Tres tests presentan \emph{flakiness} documentada en
\filepath{docs/_archive/phase1/RISKS.md}:

\begin{itemize}
    \item Proptest de equivalencia DFS$\leftrightarrow$SIMD: timeout
    intermitente con seeds que generan grafos con hubs de grado
    $>$5000.
    \item Tests del MCP contra daemon: racing en arranque si el
    daemon tarda más de lo habitual.
    \item Smoke integración con GeoIP: depende de presencia del
    archivo MaxMind; falla silenciosa si no está.
\end{itemize}

No hay fix perman\-ente aún; los dos primeros se mitigan con
retry-on-timeout en CI, el tercero con skip condicional.

\section{Evaluación}

\subsection*{Sin benchmarks comparativos formales}

No se ejecutaron benchmarks comparativos formales contra SIEMs
comerciales (Splunk, KQL, EQL). El hecho de que las superficies
sean distintas (tablas vs. grafos) hace la comparación no trivial
metodológicamente. Queda como trabajo futuro diseñar un harness
que exprese las mismas preguntas en ambos paradigmas y compare
tiempo-al-primer-resultado.

\subsection*{Datasets sintéticos de tamaño limitado}

Los smoke tests usan datasets $<$100k eventos. El comportamiento a
$10^7$ eventos (régimen productivo real) no está medido
sistemáticamente. Observaciones empíricas sobre grafos reales
sugieren que el matcher escala; no reportamos números porque no
son reproducibles.

\section{Superficies}

\subsection*{DSL sin operadores probabilísticos}

El DSL expresa patrones determinísticos (``hay una arista'').
No expresa ``hay una arista que ocurre menos del 5\% del tiempo''
o ``esta subsecuencia es inusual''. Esas preguntas se
responden via scoring post-hoc, no via el DSL. Un DSL con
cuantificadores probabilísticos es posible pero fuera de
alcance.

\subsection*{Canonical OCSF incompleto}

El mapping a OCSF cubre los campos comunes
(\code{time}, \code{actor}, \code{src\_endpoint}, \code{dst\_endpoint},
\code{observables}). No cubre todas las clases
(\code{class\_uid} de OCSF v1.4 excede lo mapeable con el
abanico actual de parsers). La extensión por clase va bajo demanda.

\subsection*{Provenance en memoria, no en disco}

El \code{provenance.mapping_hash} de un evento vive con el evento
en el grafo, pero la tabla \code{mapping_hash$\to$VRL source} se
regenera al arrancar. Si la
\code{mapping_library} se pierde, los hashes en el grafo quedan
colgando. Una persistencia a disco de la library es follow-up
conocido.

\section{Operación}

\subsection*{No hay multi-tenant}

El daemon asume un solo operador. No hay aislamiento por usuario
ni políticas de acceso a sesiones. Para deploys multi-cliente
hace falta una capa de gateway externa.

\subsection*{Sin backup automático}

La sesión persiste on-close pero no hay snapshotting periódico. Un
crash pierde el estado desde el último close.
